\documentclass[12pt,a4paper]{article}
\usepackage[utf8]{inputenc}
\usepackage[spanish]{babel}
\usepackage{geometry}
\usepackage{fancyhdr}
\usepackage{enumitem}
\usepackage{xcolor}
\usepackage{listings}
\usepackage{graphicx}
\usepackage{hyperref}
\usepackage{tabularx}
\usepackage{array}

\geometry{margin=2cm}
\pagestyle{fancy}
\fancyhf{}
\lhead{Curso de Programación en Python}
\rhead{Semana 2}
\cfoot{\thepage}

% Configuración para código Python
\lstset{
    language=Python,
    basicstyle=\ttfamily\small,
    keywordstyle=\color{blue},
    stringstyle=\color{red},
    commentstyle=\color{green!60!black},
    showstringspaces=false,
    breaklines=true,
    frame=single,
    numbers=left,
    numberstyle=\tiny\color{gray}
}

\title{\textbf{Tarea Semana 2: Juego de Adivinanzas}\\
\large Bucles y Estructuras de Repetición}
\author{Curso de Programación en Python}
\date{Tiempo de entrega: 1 semana}

\begin{document}

\maketitle

\section{Introducción}

¡Bienvenido a la Semana 2! Ya dominas los fundamentos de Python y sabes tomar decisiones con condicionales. Ahora es momento de aprender una de las herramientas más poderosas de la programación: \textbf{los bucles}.

Los bucles (loops) permiten que tu programa repita acciones múltiples veces sin tener que escribir el mismo código una y otra vez. Es como cuando preparas 10 tazas de café: no escribes 10 veces las instrucciones, simplemente repites el proceso 10 veces.

En esta tarea crearás un juego interactivo de adivinanzas que usa bucles para dar múltiples oportunidades al jugador, validar entradas, generar números aleatorios y crear una experiencia de juego completa y divertida.

\section{Objetivos de Aprendizaje}

Al completar esta tarea, serás capaz de:

\begin{itemize}[leftmargin=*]
    \item Usar bucles \texttt{while} para repetir código mientras una condición sea verdadera
    \item Usar bucles \texttt{for} para iterar sobre secuencias
    \item Trabajar con la función \texttt{range()} para generar secuencias de números
    \item Controlar el flujo de bucles con \texttt{break} y \texttt{continue}
    \item Validar entradas del usuario usando bucles
    \item Evitar bucles infinitos mediante condiciones bien diseñadas
    \item Combinar bucles con condicionales para lógica compleja
    \item Iterar sobre listas y secuencias
    \item Usar contadores y acumuladores en bucles
\end{itemize}

\section{Enunciado del Proyecto}

\subsection{Descripción General}

Vas a crear un programa llamado \textbf{``Adivinador de Números - Edición Costa Rica''}. Este es un juego interactivo donde:

\begin{itemize}
    \item El programa genera un número secreto aleatorio
    \item El jugador tiene un número limitado de intentos para adivinarlo
    \item El programa da pistas después de cada intento (muy alto, muy bajo, cerca, lejos)
    \item Se llevan estadísticas del juego
    \item El jugador puede jugar múltiples rondas
    \item Hay diferentes niveles de dificultad
\end{itemize}

El juego debe ser entretenido, dar retroalimentación útil, y demostrar tu dominio de bucles \texttt{while}, \texttt{for}, \texttt{break}, \texttt{continue}, y \texttt{range()}.

\subsection{Modos de Juego}

El programa debe ofrecer 3 niveles de dificultad:

\begin{enumerate}
    \item \textbf{Fácil:} Número entre 1 y 50, 10 intentos, pistas detalladas
    \item \textbf{Medio:} Número entre 1 and 100, 7 intentos, pistas normales
    \item \textbf{Difícil:} Número entre 1 y 200, 5 intentos, pistas básicas
\end{enumerate}

\section{Requisitos Funcionales}

\subsection{Parte 1: Menú Principal}

El programa debe iniciar con un menú que permita:

\begin{enumerate}[leftmargin=*]
    \item Jugar una nueva partida
    \item Ver estadísticas (partidas jugadas, ganadas, perdidas)
    \item Ver instrucciones del juego
    \item Salir del programa
\end{enumerate}

\textbf{Implementación requerida:}
\begin{itemize}
    \item Usar un bucle \texttt{while} para mantener el menú activo
    \item Validar que la opción elegida sea válida (1-4)
    \item Si la opción es inválida, pedir nuevamente (usar \texttt{continue})
    \item Solo salir cuando el usuario elija la opción correspondiente
\end{itemize}

\subsection{Parte 2: Selección de Dificultad}

Cuando el jugador elige jugar:

\begin{enumerate}[leftmargin=*]
    \item Mostrar los 3 niveles de dificultad con sus características
    \item Solicitar que elija un nivel (1-3)
    \item Validar la selección con un bucle
    \item Configurar el juego según el nivel elegido:
    \begin{itemize}
        \item Rango de números (máximo)
        \item Cantidad de intentos permitidos
        \item Nivel de detalle de pistas
    \end{itemize}
\end{enumerate}

\subsection{Parte 3: Mecánica del Juego}

\subsubsection{Generación del Número Secreto}

\begin{itemize}
    \item Usar \texttt{import random}
    \item Generar número aleatorio según el nivel: \texttt{random.randint(1, maximo)}
    \item El número debe mantenerse secreto durante el juego
\end{itemize}

\subsubsection{Bucle Principal del Juego}

Implementar con \texttt{while} o \texttt{for} con \texttt{range()}:

\begin{enumerate}[leftmargin=*]
    \item \textbf{Mostrar intento actual:} ``Intento 1 de 10''
    
    \item \textbf{Solicitar adivinanza del jugador:}
    \begin{itemize}
        \item Validar que sea un número
        \item Validar que esté en el rango válido
        \item Si es inválido, no contar como intento y pedir nuevamente
    \end{itemize}
    
    \item \textbf{Comparar con el número secreto:}
    \begin{itemize}
        \item Si adivinó correctamente: ¡Felicitar y terminar con \texttt{break}!
        \item Si no adivinó: Dar pistas según el nivel
    \end{itemize}
    
    \item \textbf{Sistema de pistas:}
    \begin{itemize}
        \item \textbf{Básicas (todos los niveles):} ``Muy alto'' o ``Muy bajo''
        \item \textbf{Normales (medio y fácil):} ``Cerca'' (diferencia $\leq$ 10) o ``Lejos''
        \item \textbf{Detalladas (solo fácil):} ``¡Hirviendo!'' (diferencia $\leq$ 5), ``Tibio'' (6-10), ``Frío'' (11-20), ``Congelado'' (21+)
    \end{itemize}
    
    \item \textbf{Mostrar intentos restantes:} Después de cada pista
    
    \item \textbf{Si se acaban los intentos:} Revelar el número y terminar
\end{enumerate}

\subsection{Parte 4: Historial de Intentos}

Durante el juego, mantener una lista de todos los números que el jugador ha intentado:

\begin{itemize}
    \item Usar una lista para guardar cada intento
    \item Antes de dar pistas, mostrar todos los intentos anteriores
    \item Usar un bucle \texttt{for} para mostrar el historial formateado
    \item Detectar si el jugador intenta el mismo número dos veces (dar aviso)
\end{itemize}

\subsection{Parte 5: Estadísticas Globales}

El programa debe mantener estadísticas entre partidas:

\begin{itemize}
    \item Total de partidas jugadas
    \item Partidas ganadas
    \item Partidas perdidas
    \item Mejor puntuación (menor número de intentos para ganar)
    \item Nivel más jugado
\end{itemize}

\textbf{Implementación:}
\begin{itemize}
    \item Usar variables que se actualicen después de cada partida
    \item Mostrar estadísticas cuando el usuario lo solicite del menú
    \item Usar bucles \texttt{for} para mostrar estadísticas formateadas
\end{itemize}

\subsection{Parte 6: Validación de Entradas}

Todas las entradas del usuario deben validarse con bucles:

\begin{lstlisting}[caption=Ejemplo de validación con while]
# Validar que sea un numero
while True:
    entrada = input("Ingresa tu adivinanza: ")
    if entrada.isdigit():
        numero = int(entrada)
        if 1 <= numero <= maximo:
            break  # Entrada valida, salir del bucle
        else:
            print(f"Error: Debe estar entre 1 y {maximo}")
    else:
        print("Error: Debes ingresar un numero")
\end{lstlisting}

\newpage

\subsection{Ejemplo de Ejecución Completa}

\begin{lstlisting}[caption=Flujo de juego completo]
========================================
  ADIVINADOR DE NUMEROS - COSTA RICA
========================================

--- MENU PRINCIPAL ---
1. Jugar nueva partida
2. Ver estadisticas
3. Instrucciones
4. Salir

Elige una opcion: 2

--- ESTADISTICAS ---
Partidas jugadas: 5
Partidas ganadas: 3
Partidas perdidas: 2
Mejor puntuacion: 4 intentos
Nivel mas jugado: Medio

Presiona ENTER para continuar...

--- MENU PRINCIPAL ---
1. Jugar nueva partida
2. Ver estadisticas
3. Instrucciones
4. Salir

Elige una opcion: 1

--- SELECCIONA DIFICULTAD ---
1. FACIL   - Numero entre 1 y 50  (10 intentos)
2. MEDIO   - Numero entre 1 y 100 (7 intentos)
3. DIFICIL - Numero entre 1 y 200 (5 intentos)

Elige dificultad: 2

========================================
Nivel MEDIO: Adivina el numero entre 1 y 100
Tienes 7 intentos. Buena suerte!
========================================

Intento 1 de 7
Tus intentos anteriores: (ninguno)
Ingresa tu adivinanza: 50

Muy bajo! El numero es mayor que 50.
Estas lejos del numero secreto.
Intentos restantes: 6

Intento 2 de 7
Tus intentos anteriores: [50]
Ingresa tu adivinanza: 75

Muy alto! El numero es menor que 75.
Estas cerca del numero secreto!
Intentos restantes: 5

Intento 3 de 7
Tus intentos anteriores: [50, 75]
Ingresa tu adivinanza: 65

Muy bajo! El numero es mayor que 65.
Estas cerca del numero secreto!
Intentos restantes: 4

Intento 4 de 7
Tus intentos anteriores: [50, 75, 65]
Ingresa tu adivinanza: 70

Muy bajo! El numero es mayor que 70.
Estas muy cerca del numero secreto!
Intentos restantes: 3

Intento 5 de 7
Tus intentos anteriores: [50, 75, 65, 70]
Ingresa tu adivinanza: 73

CORRECTO! El numero era 73!
Lo adivinaste en 5 intentos!

--- RESUMEN DE LA PARTIDA ---
Numero secreto: 73
Intentos usados: 5 de 7
Tus intentos: 50, 75, 65, 70, 73
Nivel: MEDIO
Resultado: GANASTE!

Quieres jugar otra vez? (s/n): s

[El juego reinicia...]
\end{lstlisting}

\newpage

\section{Requisitos Técnicos Detallados}

\subsection{Uso Obligatorio de Bucles}

Tu programa DEBE incluir:

\subsubsection{1. Bucle \texttt{while} para el menú principal}

\begin{lstlisting}
# Menu que se repite hasta que el usuario elija salir
juego_activo = True
while juego_activo:
    # Mostrar menu
    # Procesar opcion
    if opcion == 4:
        juego_activo = False  # Salir del bucle
\end{lstlisting}

\subsubsection{2. Bucle \texttt{for} con \texttt{range()} para intentos}

\begin{lstlisting}
# Dar un numero limitado de intentos
for intento in range(1, max_intentos + 1):
    print(f"Intento {intento} de {max_intentos}")
    # Pedir adivinanza
    # Si adivina correctamente:
    #     break  # Salir del bucle
\end{lstlisting}

\textbf{O usar \texttt{while} con contador:}

\begin{lstlisting}
intento = 1
while intento <= max_intentos:
    print(f"Intento {intento} de {max_intentos}")
    # Logica del juego
    intento += 1  # Incrementar contador
\end{lstlisting}

\subsubsection{3. Bucle \texttt{while True} para validación}

\begin{lstlisting}
# Validar entrada hasta que sea correcta
while True:
    try:
        numero = int(input("Ingresa un numero: "))
        if 1 <= numero <= maximo:
            break  # Entrada valida, salir
        else:
            print("Fuera de rango!")
    except ValueError:
        print("Debe ser un numero!")
\end{lstlisting}

\subsubsection{4. Bucle \texttt{for} para iterar sobre listas}

\begin{lstlisting}
# Mostrar historial de intentos
print("Tus intentos anteriores: ", end="")
for intento in lista_intentos:
    print(intento, end=" ")
print()  # Nueva linea
\end{lstlisting}

\subsubsection{5. Uso de \texttt{break} y \texttt{continue}}

\begin{lstlisting}
# break: Salir del bucle cuando se adivine
if adivinanza == numero_secreto:
    print("Correcto!")
    break

# continue: Saltar iteracion si entrada invalida
if not entrada.isdigit():
    print("Error: Solo numeros")
    continue  # Volver al inicio del bucle
\end{lstlisting}

\subsection{Estructura General del Código}

Tu código debe estar organizado así:

\begin{lstlisting}
# Encabezado con tu informacion
# Nombre, fecha, descripcion

# Importar libreria random
import random

# Variables globales para estadisticas
partidas_jugadas = 0
partidas_ganadas = 0
partidas_perdidas = 0
mejor_puntuacion = float('inf')

# Bucle principal del programa
while True:
    # Mostrar menu principal
    # Opcion 1: Jugar
    # Opcion 2: Ver estadisticas
    # Opcion 3: Instrucciones
    # Opcion 4: Salir
\end{lstlisting}

\subsection{Validaciones Requeridas}

Debes validar con bucles:

\begin{enumerate}[leftmargin=*]
    \item \textbf{Opción del menú:} Debe ser 1, 2, 3, o 4
    \item \textbf{Nivel de dificultad:} Debe ser 1, 2, o 3
    \item \textbf{Adivinanza:} Debe ser un número entero en el rango válido
    \item \textbf{Jugar otra vez:} Debe ser 's' o 'n' (aceptar mayúsculas/minúsculas)
\end{enumerate}

\section{Consejos y Recomendaciones}

\begin{itemize}[leftmargin=*]
    \item \textbf{Evita bucles infinitos:} Asegúrate que siempre hay una forma de salir del bucle
    
    \item \textbf{Usa \texttt{break} sabiamente:} Para salir de bucles cuando se cumpla una condición
    
    \item \textbf{Usa \texttt{continue} para saltar:} Cuando una iteración no deba completarse
    
    \item \textbf{Incrementa contadores:} Si usas \texttt{while}, no olvides incrementar el contador
    
    \item \textbf{Prueba diferentes niveles:} Asegúrate que cada dificultad funciona correctamente
    
    \item \textbf{Prueba casos extremos:}
    \begin{itemize}
        \item Adivinar en el primer intento
        \item Agotar todos los intentos
        \item Ingresar el número máximo
        \item Ingresar el número mínimo (1)
        \item Ingresar valores inválidos
    \end{itemize}
    
    \item \textbf{Usa variables descriptivas:}
    \begin{itemize}
        \item \texttt{numero\_secreto}, no \texttt{x}
        \item \texttt{intento\_actual}, no \texttt{i}
        \item \texttt{max\_intentos}, no \texttt{n}
    \end{itemize}
    
    \item \textbf{Comenta bucles complejos:} Explica qué hace cada bucle y por qué
\end{itemize}

\section{Entregables}

\begin{enumerate}[leftmargin=*]
    \item Archivo Python llamado \texttt{adivinador\_numeros.py}
    \item El código debe ejecutarse sin errores
    \item Debe incluir todos los comentarios requeridos
    \item Debe implementar todos los niveles de dificultad
    \item Debe mantener estadísticas entre partidas
\end{enumerate}

\section{Defensa del Código (60\% de la nota final)}

\subsection{¿Qué es la defensa del código?}

La defensa del código es una reunión virtual donde presentarás y explicarás tu trabajo. Esta defensa es \textbf{obligatoria} y representa el \textbf{60\% de la nota final} de esta tarea.

\subsection{¿Cómo funciona?}

Durante la defensa deberás:

\begin{enumerate}[leftmargin=*]
    \item \textbf{Compartir tu pantalla} mostrando tu código en el editor.
    
    \item \textbf{Ejecutar tu programa} demostrando:
    \begin{itemize}
        \item El menú principal funcionando
        \item Una partida ganada
        \item Una partida perdida
        \item Las estadísticas actualizándose
        \item Validación de entradas inválidas
        \item Los 3 niveles de dificultad
    \end{itemize}
    
    \item \textbf{Explicar tus bucles:}
    \begin{itemize}
        \item Por qué elegiste \texttt{while} vs \texttt{for} en cada caso
        \item Cómo funciona tu bucle de validación
        \item Qué condición hace que cada bucle termine
        \item Cómo usas \texttt{break} y \texttt{continue}
        \item Cómo evitaste bucles infinitos
    \end{itemize}
    
    \item \textbf{Responder preguntas como:}
    \begin{itemize}
        \item ``¿Qué pasa si quitas este \texttt{break}?''
        \item ``¿Por qué usaste \texttt{range(1, max\_intentos + 1)} y no \texttt{range(max\_intentos)}?''
        \item ``¿Cómo funciona el sistema de pistas?''
        \item ``¿Qué hace \texttt{continue} en este bucle?''
        \item ``¿Cómo mantienes las estadísticas entre partidas?''
        \item ``¿Qué pasaría si el usuario nunca elige salir?''
    \end{itemize}
    
    \item \textbf{Demostrar comprensión de:}
    \begin{itemize}
        \item La diferencia entre \texttt{while} y \texttt{for}
        \item Cuándo usar cada tipo de bucle
        \item Cómo funcionan \texttt{break} y \texttt{continue}
        \item Qué es un bucle infinito y cómo evitarlo
        \item Cómo iterar sobre listas
    \end{itemize}
\end{enumerate}

\subsection{Preguntas Típicas de la Defensa}

Prepárate para responder:

\begin{itemize}[leftmargin=*]
    \item ¿Cuál es la diferencia entre \texttt{while} y \texttt{for}?
    \item ¿Por qué elegiste \texttt{while True} para la validación?
    \item ¿Qué hace \texttt{range(1, 11)}? ¿Incluye el 11?
    \item ¿Cómo funciona el \texttt{break}? Dame un ejemplo de tu código
    \item ¿Qué pasaría si olvidas incrementar el contador en un \texttt{while}?
    \item ¿Cómo iteraste sobre la lista de intentos?
    \item Si cambio \texttt{break} por \texttt{continue}, ¿qué pasa?
    \item ¿Cómo garantizas que el menú siempre vuelva a mostrarse?
    \item ¿Por qué es importante validar antes de procesar?
\end{itemize}

\subsection{Consejos para una Buena Defensa}

\begin{itemize}[leftmargin=*]
    \item \textbf{Conoce tus bucles:} Entiende perfectamente cada \texttt{while} y \texttt{for} que usaste
    \item \textbf{Practica modificaciones:} ¿Qué pasa si cambias una condición?
    \item \textbf{Explica el flujo:} Describe cómo el programa pasa por los bucles
    \item \textbf{Ten ejemplos listos:} Prepara casos de prueba interesantes
    \item \textbf{Entiende \texttt{break} y \texttt{continue}:} Son conceptos clave de esta semana
\end{itemize}

\subsection{Distribución de la Nota}

\begin{itemize}
    \item \textbf{Código funcional (40\%):} El programa funciona según los requisitos
    \item \textbf{Defensa del código (60\%):} Tu capacidad de explicar y defender tu trabajo
\end{itemize}

\textbf{Importante:} Si no asistes a la defensa o no puedes explicar tus bucles, la nota máxima será 40/100.

\section{Desafíos Opcionales (Puntos Extra)}

\begin{enumerate}[leftmargin=*]
    \item \textbf{Modo multijugador (+6 pts):}
    \begin{itemize}
        \item 2 jugadores se turnan para adivinar
        \item Gana quien adivine en menos intentos
        \item Mantener estadísticas por jugador
    \end{itemize}
    
    \item \textbf{Pistas avanzadas (+5 pts):}
    \begin{itemize}
        \item ``El número es par/impar''
        \item ``Es múltiplo de X''
        \item ``Es primo'' (si has investigado números primos)
    \end{itemize}
    
    \item \textbf{Sistema de puntuación (+5 pts):}
    \begin{itemize}
        \item Calcular puntos: (intentos\_restantes $\times$ 10) + bonus\_dificultad
        \item Mantener tabla de mejores puntuaciones (top 5)
        \item Mostrar ranking al finalizar
    \end{itemize}
    
    \item \textbf{Modo contrarreloj (+6 pts):}
    \begin{itemize}
        \item Agregar límite de tiempo usando \texttt{time}
        \item Mostrar tiempo restante
        \item Terminar si se acaba el tiempo
    \end{itemize}
    
    \item \textbf{Generador de patrones (+4 pts):}
    \begin{itemize}
        \item Además del juego, un modo que genera patrones numéricos con \texttt{for}
        \item Pirámides, secuencias, tablas de multiplicar
        \item Al menos 3 patrones diferentes
    \end{itemize}
\end{enumerate}

\textit{Nota: Los puntos extra solo se otorgan si la implementación es correcta Y puedes explicarla en la defensa.}

\newpage

\section{Rúbrica de Evaluación}

\subsection{Distribución General}

\begin{itemize}
    \item \textbf{Código Funcional:} 40\% (40 puntos)
    \item \textbf{Defensa del Código:} 60\% (60 puntos)
    \item \textbf{Total:} 100 puntos
    \item \textbf{Puntos Extra:} Hasta +26 puntos adicionales
\end{itemize}

\subsection{Parte 1: Código Funcional (40 puntos)}

\subsubsection{1.1 Implementación de Bucles (18 puntos)}

\begin{itemize}[leftmargin=*]
    \item \textbf{Bucle \texttt{while} para menú principal (4 pts):}
    \begin{itemize}
        \item 4 pts: Menú funcional que se repite hasta salir
        \item 3 pts: Funciona pero con pequeños errores
        \item 1 pt: Implementación parcial o incorrecta
        \item 0 pts: No implementa menú con bucle
    \end{itemize}
    
    \item \textbf{Bucle para intentos del juego (5 pts):}
    \begin{itemize}
        \item 5 pts: Usa \texttt{for} con \texttt{range()} o \texttt{while} correctamente
        \item 4 pts: Funciona pero lógica mejorable
        \item 2 pts: Funciona parcialmente
        \item 0 pts: No implementa bucle de intentos
    \end{itemize}
    
    \item \textbf{Bucles de validación (5 pts):}
    \begin{itemize}
        \item 5 pts: Valida todas las entradas con bucles apropiados
        \item 4 pts: Valida la mayoría de entradas
        \item 2 pts: Validación mínima
        \item 0 pts: No valida entradas con bucles
    \end{itemize}
    
    \item \textbf{Uso de \texttt{break} y \texttt{continue} (4 pts):}
    \begin{itemize}
        \item 4 pts: Usa ambos correctamente en contextos apropiados
        \item 3 pts: Usa uno correctamente
        \item 1 pt: Intenta usar pero con errores
        \item 0 pts: No usa \texttt{break} ni \texttt{continue}
    \end{itemize}
\end{itemize}

\subsubsection{1.2 Funcionalidad del Juego (12 puntos)}

\begin{itemize}[leftmargin=*]
    \item \textbf{Generación de número aleatorio (3 pts):}
    \begin{itemize}
        \item 3 pts: Genera correctamente según nivel de dificultad
        \item 2 pts: Genera pero con errores menores
        \item 1 pt: Implementación parcial
        \item 0 pts: No genera número aleatorio
    \end{itemize}
    
    \item \textbf{Sistema de pistas (4 pts):}
    \begin{itemize}
        \item 4 pts: Pistas precisas y diferenciadas por nivel
        \item 3 pts: Pistas funcionales pero no diferenciadas
        \item 2 pts: Pistas básicas solamente
        \item 0 pts: No da pistas
    \end{itemize}
    
    \item \textbf{Niveles de dificultad (3 pts):}
    \begin{itemize}
        \item 3 pts: 3 niveles funcionando correctamente
        \item 2 pts: 2 niveles funcionando
        \item 1 pt: Solo 1 nivel
        \item 0 pts: No implementa niveles
    \end{itemize}
    
    \item \textbf{Detección de victoria/derrota (2 pts):}
    \begin{itemize}
        \item 2 pts: Detecta ambos casos correctamente y termina apropiadamente
        \item 1 pt: Detecta pero con errores
        \item 0 pts: No detecta correctamente
    \end{itemize}
\end{itemize}

\subsubsection{1.3 Características Adicionales (10 puntos)}

\begin{itemize}[leftmargin=*]
    \item \textbf{Historial de intentos (4 pts):}
    \begin{itemize}
        \item 4 pts: Lista completa con iteración usando \texttt{for}
        \item 3 pts: Mantiene lista pero presentación mejorable
        \item 1 pt: Implementación parcial
        \item 0 pts: No mantiene historial
    \end{itemize}
    
    \item \textbf{Estadísticas globales (4 pts):}
    \begin{itemize}
        \item 4 pts: Mantiene todas las estadísticas entre partidas
        \item 3 pts: Mantiene algunas estadísticas
        \item 1 pt: Estadísticas básicas
        \item 0 pts: No mantiene estadísticas
    \end{itemize}
    
    \item \textbf{Opción de jugar múltiples partidas (2 pts):}
    \begin{itemize}
        \item 2 pts: Permite jugar múltiples veces sin reiniciar programa
        \item 1 pt: Funciona con errores
        \item 0 pts: Requiere reiniciar programa
    \end{itemize}
\end{itemize}

\subsection{Parte 2: Defensa del Código (60 puntos)}

\subsubsection{2.1 Comprensión de Bucles (30 puntos)}

\begin{itemize}[leftmargin=*]
    \item \textbf{Explicación de \texttt{while} vs \texttt{for} (10 pts):}
    \begin{itemize}
        \item 10 pts: Explica perfectamente diferencias y cuándo usar cada uno
        \item 7 pts: Explica correctamente con pequeñas imprecisiones
        \item 4 pts: Explicación básica o parcial
        \item 0 pts: No entiende la diferencia
    \end{itemize}
    
    \item \textbf{Explicación de \texttt{break} y \texttt{continue} (8 pts):}
    \begin{itemize}
        \item 8 pts: Explica claramente ambos con ejemplos de su código
        \item 6 pts: Explica correctamente con alguna duda
        \item 3 pts: Explicación superficial
        \item 0 pts: No entiende su funcionamiento
    \end{itemize}
    
    \item \textbf{Explicación de \texttt{range()} (6 pts):}
    \begin{itemize}
        \item 6 pts: Explica perfectamente cómo funciona y sus parámetros
        \item 4 pts: Explica correctamente lo básico
        \item 2 pts: Explicación confusa
        \item 0 pts: No entiende \texttt{range()}
    \end{itemize}
    
    \item \textbf{Prevención de bucles infinitos (6 pts):}
    \begin{itemize}
        \item 6 pts: Explica cómo garantizó que cada bucle termina
        \item 4 pts: Comprensión básica del concepto
        \item 2 pts: Explicación vaga
        \item 0 pts: No comprende el concepto
    \end{itemize}
\end{itemize}

\subsubsection{2.2 Demostración y Razonamiento (20 puntos)}

\begin{itemize}[leftmargin=*]
    \item \textbf{Ejecución de diferentes escenarios (8 pts):}
    \begin{itemize}
        \item 8 pts: Demuestra exitosamente múltiples escenarios (ganar, perder, validaciones)
        \item 6 pts: Demuestra varios escenarios
        \item 3 pts: Demuestra escenarios básicos
        \item 0 pts: No puede demostrar diferentes casos
    \end{itemize}
    
    \item \textbf{Respuesta a preguntas técnicas (12 pts):}
    \begin{itemize}
        \item 12 pts: Responde correcta y claramente todas las preguntas sobre bucles
        \item 9 pts: Responde bien la mayoría
        \item 5 pts: Respuestas parciales
        \item 2 pts: Dificultad con preguntas básicas
        \item 0 pts: No puede responder sobre su lógica
    \end{itemize}
\end{itemize}

\subsubsection{2.3 Calidad del Código (10 puntos)}

\begin{itemize}[leftmargin=*]
    \item \textbf{Comentarios explicativos (4 pts):}
    \begin{itemize}
        \item 4 pts: Comentarios claros explicando bucles complejos
        \item 3 pts: Comentarios presentes pero básicos
        \item 1 pt: Comentarios mínimos
        \item 0 pts: Sin comentarios
    \end{itemize}
    
    \item \textbf{Nombres de variables (3 pts):}
    \begin{itemize}
        \item 3 pts: Todos los nombres son descriptivos
        \item 2 pts: Mayoría de nombres buenos
        \item 1 pt: Nombres poco claros
        \item 0 pts: Nombres inapropiados
    \end{itemize}
    
    \item \textbf{Organización y legibilidad (3 pts):}
    \begin{itemize}
        \item 3 pts: Código perfectamente organizado y legible
        \item 2 pts: Código legible con errores menores
        \item 1 pt: Poco organizado
        \item 0 pts: Código desorganizado
    \end{itemize}
\end{itemize}

\subsection{Puntos Extra (Hasta +26 puntos)}

\begin{itemize}[leftmargin=*]
    \item \textbf{Modo multijugador (+6 pts):} Implementa sistema de turnos entre 2 jugadores
    \item \textbf{Pistas avanzadas (+5 pts):} Pistas matemáticas (par/impar, múltiplos, primos)
    \item \textbf{Sistema de puntuación (+5 pts):} Calcula puntos y mantiene top 5
    \item \textbf{Modo contrarreloj (+6 pts):} Límite de tiempo con módulo \texttt{time}
    \item \textbf{Generador de patrones (+4 pts):} Genera 3+ patrones numéricos con bucles
\end{itemize}

\textbf{Nota:} Los puntos extra solo se otorgan si la implementación es correcta Y el estudiante puede explicarla completamente durante la defensa.

\subsection{Penalizaciones}

\begin{itemize}[leftmargin=*]
    \item \textbf{-100\% de la defensa:} No asiste sin justificación (nota máxima = 40/100)
    \item \textbf{-50\% de la defensa:} Evidencia de no entender el código
    \item \textbf{-5 pts:} Entrega tardía de 1 día
    \item \textbf{-10 pts:} Entrega tardía de 2-3 días
    \item \textbf{-20 pts:} Entrega tardía de 4+ días
    \item \textbf{-3 pts:} Nombre de archivo incorrecto
    \item \textbf{-5 pts:} Bucle infinito en el código (no termina nunca)
    \item \textbf{-3 pts:} No usa \texttt{import random} correctamente
\end{itemize}

\subsection{Nota Mínima Aprobatoria}

Para aprobar esta tarea necesitas obtener al menos \textbf{60 puntos de 100}.

\subsection{Escala de Calificación}

\begin{itemize}
    \item \textbf{90-100+ pts:} Excelente - Dominio completo de bucles
    \item \textbf{80-89 pts:} Muy Bueno - Comprensión sólida de estructuras de repetición
    \item \textbf{70-79 pts:} Bueno - Comprensión aceptable, necesita práctica
    \item \textbf{60-69 pts:} Suficiente - Comprensión básica, necesita refuerzo
    \item \textbf{0-59 pts:} Insuficiente - No aprobado
\end{itemize}

\section{Recursos de Apoyo}

\begin{itemize}[leftmargin=*]
    \item \textbf{Documentación oficial:} \url{https://docs.python.org/es/3/tutorial/controlflow.html}
    \item \textbf{Módulo random:} \url{https://docs.python.org/es/3/library/random.html}
    \item \textbf{Bucles while:} Revisa las notas sobre condiciones de terminación
    \item \textbf{Bucles for:} Practica con \texttt{range()} en el intérprete interactivo
    \item \textbf{Break y continue:} Experimenta con ejemplos simples primero
\end{itemize}

\section{Preguntas Frecuentes}

\textbf{P: ¿Cuál es la diferencia entre \texttt{while} y \texttt{for}?}\\
R: \texttt{while} repite mientras una condición sea verdadera. \texttt{for} repite un número específico de veces o sobre una secuencia.

\textbf{P: ¿Qué hace \texttt{range(1, 11)}?}\\
R: Genera números del 1 al 10 (el límite superior 11 NO se incluye).

\textbf{P: ¿Cómo evito un bucle infinito?}\\
R: Asegúrate que la condición del \texttt{while} eventualmente sea falsa, o que haya un \texttt{break} en algún punto.

\textbf{P: ¿Cuándo uso \texttt{break} vs \texttt{continue}?}\\
R: \texttt{break} sale completamente del bucle. \texttt{continue} salta a la siguiente iteración.

\textbf{P: ¿Cómo genero un número aleatorio?}\\
R: \texttt{import random} y luego \texttt{random.randint(1, 100)} para un número entre 1 y 100.

\textbf{P: ¿Puedo usar listas?}\\
R: Sí! Para el historial de intentos es perfecto usar una lista.

\textbf{P: ¿Es obligatorio hacer los 3 niveles de dificultad?}\\
R: Sí, es un requisito funcional obligatorio.

\vspace{1cm}

\begin{center}
\textbf{¡Éxito en tu juego de adivinanzas!}\\
\textit{Los bucles son el corazón de la programación.\\
Domínalos y podrás automatizar cualquier tarea repetitiva.}
\end{center}

\end{document}
